%=====================================================================
\section{Ejercicios adicionales}
%=====================================================================

% =====================================================================
\subsection{Ejercicio --- Bingo virtual}
% =====================================================================

\begin{consigna}
Una franquicia de bingos decide crear una aplicación que simule un bingo virtual donde
los participantes paguen y obtengan dinero real.

Cada juego tiene un valor de ingreso, y le permite al participante obtener un cartón con
15 número elegidos al azar entre 1 y 90. Cuando comienza el juego, el servidor comienza a
``sacar'' números, y los participantes deben ir marcando los números en su cartón si
corresponde.

Cuando uno de los participantes completa los 15 número de su cartón, selecciona el botón
``bingo'', y si realmente salieron todos, gana el premio de esa jugada y termina la
partida para todos. Si no fuese así, se le avisa y se lo descalifica.

Resuelta la seguridad de los aspectos de introducir y retirar dinero se identificaron
tres funcionalidades necesarias respecto a la dinámica del juego. Explicar como se podría
lograr cada una recurriendo a funciones criptográficas \underline{simétricas solamente},
justificando la elección. De ser necesario, explicar que información se intercambia en
cada momento entre las apps y el servidor.

\begin{enumerate}
  \item[a)] El servidor debe generar una secuencia de números entre 1 y 90, que luego
  pueda ser reproducida por los clientes al finalizar la partida (para garantizar que no
  se favoreció adrede a alguien en el sorteo). (1 pto.)
  \item[b)] Los cartones se crean en las apps cliente al inicio de la partida. El servidor
  no conoce los cartones hasta el momento de verificar un bingo, pero debe ser posible
  para el servidor verificar que el cartón no fue modificado una vez empezado el juego
  (para evitar que un jugador arme el cartón ideal para ganar luego de 15 números). (1 pto.)
  \item[c)] La app de cada participante deberá poder verificar si el ganador realmente
  completó el bingo y emitir una alerta en caso contrario (para evitar que el servidor de
  por ganador a alguien que no completo el cartón). Las apps no se comunican entre si,
  sino a través del servidor. (1 pto.)
\end{enumerate}

\emph{Nota}: Considerar que hay aproximadamente $2^{55}$ cartones (combinatoria de 90
tomados de a 15), y asumir que un atacante podría llegar a probar esa cantidad de
combinaciones.
\end{consigna}

\paragraph{Temas.}
Clase 3 --- MACs y Cifrado Autenticado (\S Qué es un MAC; \S HMAC; \S funciones de hash y
sus propiedades de resistencia; \S usos de la resistencia a preimagen para protocolos de
compromiso/prueba de conocimiento; \S no-repudio: por qué los MAC no lo dan).

\paragraph{Qué necesitás saber.}
\begin{itemize}
  \item \textbf{MAC/HMAC:} terna $(\mathrm{Gen},\mathrm{Mac},\mathrm{Vrfy})$ con clave
  \textbf{compartida} $k$; da integridad y autenticación de origen \emph{entre partes que
  comparten la clave}, pero \textbf{no} no-repudio.
  \item \textbf{Función de hash:} sin clave, de tamaño fijo, con tres resistencias
  (preimagen, segunda preimagen, colisión). La resistencia a preimagen es la base de los
  \textbf{protocolos de prueba de conocimiento / compromiso}: publicar $H(x)$ y después
  revelar $x$, de forma que nadie pudo deducir $x$ antes de la revelación, y quien reveló
  $x$ no puede haber cambiado de valor sin que se note (segunda preimagen).
  \item \textbf{Ataque de cumpleaños / brute-force de espacios chicos:} si el espacio de
  posibles mensajes es chico (acá, $2^{55}$ cartones), un compromiso $H(\text{cartón})$
  \emph{sin sal} es vulnerable a que cualquiera enumere los $2^{55}$ cartones posibles,
  los hashee y compare contra el compromiso publicado para deducir el cartón antes de
  tiempo. La nota del enunciado es justamente el gancho para exigir un \textbf{nonce/sal
  grande} (mucho mayor a 55 bits) dentro del compromiso.
\end{itemize}

\paragraph{Notas y conexiones.}
Los tres incisos son variaciones del mismo patrón: (a) y (b) son \textbf{compromisos}
(commitments) usando la resistencia a preimagen/segunda preimagen de una función de hash
(no llevan clave, pero son ``funciones criptográficas simétricas'' en el sentido amplio
que exige la consigna, ya que no dependen de un par de claves pública/privada); (c) es un
esquema de \textbf{autenticación con clave compartida} (MAC/HMAC) para que todas las apps
puedan validar de forma independiente los mismos mensajes del servidor. Como está prohibido
usar asimétrica, no hay firma digital disponible, así que ninguna de las tres soluciones
da no-repudio frente al servidor (el servidor y las apps comparten secretos).

\paragraph{Resolución.}

\textbf{a) Sorteo verificable.} Antes de arrancar la partida el servidor elige una
\textbf{semilla secreta} $s$ (número aleatorio, e.g. 256 bits) y publica un
\textbf{compromiso} $c=H(s)$ a todas las apps (usando la resistencia a preimagen de un
hash: nadie puede deducir $s$ a partir de $c$). La secuencia de 90 números se deriva de
forma \textbf{determinística} de $s$ con una función pseudoaleatoria con clave (p.ej.
$\mathrm{HMAC}_s(i)$ para $i=1,2,\dots$, usada como fuente de aleatoriedad de un
\emph{shuffle} tipo Fisher--Yates sobre $\{1,\dots,90\}$). El servidor va revelando los
números uno a uno a medida que los ``saca''; como se derivan de $s$ que solo él conoce, no
puede ser forzado a revelar otro número, pero tampoco puede cambiar el orden a mitad de
partida sin que se note (porque ya se comprometió a $c=H(s)$ antes de empezar). Al
finalizar la partida, el servidor revela $s$; cada app recomputa $H(s)$ y verifica que
coincide con $c$ (garantiza que $s$ es el mismo que se comprometió al inicio, por
resistencia a segunda preimagen), y recomputa la secuencia completa de 90 números con $s$
para verificar que coincide con la secuencia efectivamente anunciada durante el juego. Si
alguna etapa no coincide, la secuencia fue manipulada.

\textbf{b) Cartón no modificable.} Al empezar la partida, cada app cliente genera su
cartón y envía al servidor un \textbf{compromiso} $h=H(\text{cartón}\,\|\,r)$, donde $r$ es
un \textbf{nonce aleatorio grande} (mucho mayor a 55 bits, p.ej. 128 bits) generado por la
app y que \textbf{no} se envía al servidor en ese momento. Por la nota del enunciado, si se
omitiera $r$ el compromiso $H(\text{cartón})$ sería vulnerable a que cualquiera (el propio
servidor u otro atacante con el compromiso) recorra los $2^{55}$ cartones posibles,
hasheándolos hasta encontrar el que coincide, revelando el cartón antes de que termine la
partida. Agregar $r$ (con un espacio mucho mayor que $2^{55}$) hace que la búsqueda por
fuerza bruta sea inviable ($2^{55+128}$ combinaciones), preservando la \emph{ocultación}
del cartón hasta el final. Al momento de reclamar bingo (o al finalizar la partida), la app
envía $(\text{cartón}, r)$ al servidor, que recalcula $H(\text{cartón}\,\|\,r)$ y verifica
que coincide con el $h$ recibido al inicio (resistencia a segunda preimagen: no puede
existir otro cartón$'$ con el mismo $h$, así que el jugador no pudo cambiar el cartón
después de comprometerse).

\textbf{c) Verificación distribuida sin comunicación entre apps.} Todas las apps
comparten con el servidor (instalada de antemano, al registrarse) una \textbf{clave
simétrica} $k$ (puede ser una por partida, distribuida al inicio junto con $(G,q,g)$... o
más simple, generada por el servidor y enviada cifrada/instalada a cada app al arrancar la
partida). Cada número que el servidor ``saca'' en el paso (a) se transmite junto con una
etiqueta $t_i=\mathrm{HMAC}_k(i\,\|\,\text{número}_i)$; cada app, al recibirlo, valida
$\mathrm{Vrfy}_k$ y solo si es válido lo marca como ``anunciado''. De esta forma todas las
apps arman localmente el mismo historial \textbf{autenticado} de números sacados, sin
hablar entre sí (solo escuchan al servidor). Cuando alguien grita bingo, el servidor
retransmite a todas las apps el cartón revelado del ganador junto con $(r,h)$ del punto
(b); cada app, de forma independiente, (i) recalcula $H(\text{cartón}\|r)$ y compara con el
$h$ que el servidor había repartido al inicio de la partida (si no coincide, el cartón fue
alterado), y (ii) chequea que los 15 números del cartón estén todos entre los que su propio
historial local marcó como ``anunciados y autenticados con $k$''. Si cualquiera de las dos
verificaciones falla, la app emite una alerta al servidor (sigue sin comunicarse
directamente con las demás apps, tal como pide la consigna).

\begin{ojo}
Ninguna de las tres soluciones da \textbf{no-repudio}: como todo se basa en hash sin clave
o en un secreto \textbf{compartido} ($s$, $r$, $k$), el servidor sigue siendo, en última
instancia, quien controla la información y a quien hay que confiarle la ejecución honesta
del protocolo (el enunciado no permite usar firma digital, que sería la única forma de dar
no-repudio, según Clase 4).
\end{ojo}

% =====================================================================
\subsection{Ejercicio --- Robo de cables eléctricos}
% =====================================================================

\begin{consigna}
El robo de cables por el valor intrínseco del cobre que traen es un gran problema para las
compañías de distribución eléctrica, dado que les genera un doble perjuicio: el costo
económico de reponer los cables y el costo asociado a resarcir a clientes cuyo servicio se
ve afectado.

Para detectar cortes (intencionales o no) en las lineas, cada cierta distancia se conectan
pequeños dispositivos que modulan una señal de alta frecuencia en la misma linea que
transporta energía. Cada dispositivo tiene una firma única (una frecuencia y un periodo
dado --- diseñado de tal forma que todos los dispositivos pueden inyectar su señal sin
interferirse entre si). Las lineas terminan en transformadores. Antes del transformador se
coloca un dispositivo llamado \emph{colector} que recolecta todas las señales y aplica un
filtro de alta frecuencia que limpia la señal para que llegue al transformador sin este
ruido.

Por la forma en que se juntan las señales, cada dispositivo es capaz de transmitir el
equivalente a una secuencia de 320 bits arbitraria.

Una implementación en particular de este sistema asigna a cada dispositivo un par de claves
publica y privada de 320 bits, y hace que cada dispositivo envíe regularmente la firma
digital de su numero de serie único (solo la firma!).

El \emph{colector} conoce los números de serie de todos los dispositivos, y de cada uno la
frecuencia (lo que le permite extraer los 320 bits de la señal especifica), la ubicación
relativa en el cable y su clave publica. De esta manera, puede determinar que si deja de
recibir la señal del dispositivo $x$ y todos los anteriores hubo un corte en el cable,
mientras que si solo falta $x$ hubo un problema con el dispositivo en particular.

\begin{enumerate}
  \item[a)] Explicar paso a paso que tiene que hacer el colector para validar el estado de
  la linea y llegar a uno de tres estados: todo ok, corte cercano al dispositivo $x$,
  problemas con el dispositivo $x$. Ser preciso en la aplicación de cada función
  criptográfica involucrada y sus parámetros. (1 pto)
  \item[b)] En busca de abaratar costos, se probó una modificación al sistema donde
  \underline{todos} los dispositivo comparten una \underline{única} clave secreta con el
  aparato terminal (aunque cada uno sigue teniendo un numero de serie diferente), y
  transmiten una etiqueta de su numero de serie, generada con un MAC criptográficamente
  seguro. ¿Como cambia el algoritmo del colector? ¿Que desventajas genera este cambio a
  nivel seguridad? (1 pto)
  \item[b)] ¿Como modificaría el protocolo para evitar un ataque donde alguien corte una
  sección pero agregue un aparato que repita las señales enviadas previamente por los
  dispositivos que quedaron fuera de linea? (1 pto)
\end{enumerate}

\emph{Nota}: es imposible modificar el ancho de banda de cada dispositivo de 320 bits.
\end{consigna}

\paragraph{Temas.}
Clase 4 --- Cifrado Asimétrico y Firma Digital (\S Firmas digitales; \S RSA-Signature y por
qué se firma el hash; \S curvas elípticas: 320 bits $\approx$ seguridad de RSA-2048);
Clase 3 --- MACs y Cifrado Autenticado (\S Qué es un MAC; \S no-repudio: por qué los MAC no
lo dan); Clase 5 --- Protocolos (\S nonce y timestamp como mecanismos de frescura;
\S ataques de replay y sus defensas).

\paragraph{Qué necesitás saber.}
\begin{itemize}
  \item \textbf{Firma digital:} $\mathrm{Sign}_{sk}(m)$ con la clave \textbf{privada} del
  firmante; $\mathrm{Vrfy}_{pk}(m,\sigma)$ con la clave \textbf{pública}, que cualquiera
  (incluido el colector) puede ejecutar sin compartir secretos. Da no-repudio e integridad,
  y es \textbf{públicamente verificable}: el colector solo necesita la clave pública de
  cada dispositivo, no un secreto compartido.
  \item \textbf{320 bits como firma sobre curva elíptica:} el material de la clase indica
  que para una seguridad equivalente a RSA-2048 ($\sim$128 bits) alcanza con
  $\sim$320 bits en curvas elípticas (ECDSA/ElGamal sobre curva). Esto encaja exactamente
  con el ancho de banda de 320 bits que puede transmitir cada dispositivo: es coherente con
  que la ``firma digital de su número de serie'' quepa en una única transmisión.
  \item \textbf{MAC con clave compartida entre todos:} da integridad/autenticación pero
  \textbf{no} distingue quién generó la etiqueta (cualquiera con la clave puede forjar la
  de cualquier otro) y \textbf{no} da no-repudio (Clase 3).
  \item \textbf{Frescura (nonce/timestamp) contra replay:} Clase 5 muestra que sin un
  mecanismo de frescura, un atacante puede grabar y reenviar (\emph{replay}) un mensaje
  válido viejo. La defensa es incluir un \textbf{nonce} o \textbf{timestamp} en lo que se
  firma/etiqueta, de forma que el verificador rechace mensajes ``viejos''.
\end{itemize}

\paragraph{Notas y conexiones.}
El enunciado arma un caso de uso real de firma digital asimétrica (Clase 4) con la
restricción de ancho de banda de 320 bits, que \textbf{calza justo} con el tamaño de firma
de un esquema sobre curva elíptica visto en la Clase 4. La parte (b) reproduce el
contraste MAC vs.\ firma digital de la Clase 3 (no-repudio) y la parte (b) repetida conecta
con los mecanismos de frescura/anti-replay de la Clase 5 (Needham--Schroeder,
Denning--Sacco), pero con la restricción extra de que el tamaño transmitido no puede
crecer: se resuelve firmando/etiquetando el \textbf{hash} de (serie $\|$ contador), tal
como en Clase 4 se explica por qué se firma el hash y no el mensaje (permite mensajes de
tamaño arbitrario con una firma de tamaño fijo).

\paragraph{Resolución.}

\textbf{a) Algoritmo del colector (esquema original, con firma digital).} El colector
mantiene, para cada dispositivo $i=1,\dots,n$ ordenado por \textbf{ubicación relativa en el
cable} (desde el más cercano al colector hasta el más lejano), su número de serie
$\text{serie}_i$, su frecuencia/periodo (para poder sintonizar y extraer los 320 bits que
le corresponden) y su clave pública $pk_i$. Paso a paso:
\begin{enumerate}
  \item Para cada dispositivo $i$, el colector \textbf{sintoniza} su frecuencia/periodo
  específica y trata de capturar los 320 bits transmitidos en esa señal.
  \item Si capta señal en la frecuencia de $i$: interpreta esos 320 bits como una firma
  $\sigma_i$ y ejecuta $\mathrm{Vrfy}_{pk_i}(\text{serie}_i,\sigma_i)$. Si la verificación
  da $1$, el dispositivo $i$ está \textbf{transmitiendo correctamente} (nadie puede haber
  falsificado esa firma sin la clave privada de $i$).
  \item Si \textbf{no} capta señal en la frecuencia de un dispositivo $x$ y tampoco en la
  de \textbf{todos} los dispositivos ubicados a partir de $x$ hacia el otro extremo del
  cable (los que están ``detrás'' de $x$ respecto del colector): como la señal viaja por
  el mismo conductor de cobre, un corte físico en el punto donde está $x$ interrumpe la
  continuidad eléctrica para todos los dispositivos más allá del corte. Estado:
  \textbf{corte cercano al dispositivo $x$}.
  \item Si \textbf{solo falta} la señal de $x$, pero los dispositivos ubicados antes
  \textbf{y} después de $x$ en el cable sí llegan (y verifican $\mathrm{Vrfy}$ = 1): el
  cable sigue eléctricamente continuo (si no, tampoco llegarían los de más allá de $x$),
  así que la falla es puntual del propio dispositivo $x$ (rotura, batería, etc.). Estado:
  \textbf{problema con el dispositivo $x$}.
  \item Si todos los dispositivos verifican correctamente: estado \textbf{todo ok}.
\end{enumerate}

\textbf{b) Variante con clave secreta única (MAC compartido).} El algoritmo del colector
cambia únicamente en el paso de verificación: en vez de $\mathrm{Vrfy}_{pk_i}$ con una
clave pública distinta por dispositivo, el colector (que ahora debe conocer la
\textbf{única} clave secreta $K$ compartida por todos los dispositivos) recibe la etiqueta
$t_i=\mathrm{MAC}_K(\text{serie}_i)$ transmitida en los 320 bits, recalcula
$\mathrm{MAC}_K(\text{serie}_i)$ y compara con $t_i$ (en vez de verificar una firma con
clave pública, recalcula y compara, como cualquier MAC). El resto de la lógica (corte
cercano a $x$ vs.\ problema puntual en $x$, según qué señales llegan) no cambia.

\textbf{Desventajas de seguridad:}
\begin{itemize}
  \item \textbf{Punto único de falla:} al ser una clave \textbf{compartida por todos los
  dispositivos}, si un atacante extrae físicamente $K$ de \emph{un solo} dispositivo (por
  ejemplo abriéndolo), \textbf{compromete a todo el sistema}: puede generar etiquetas
  válidas para cualquier número de serie, falsificando la señal de cualquier dispositivo
  (incluso de los que nunca tocó).
  \item \textbf{Sin no-repudio ni atribución:} como todos comparten $K$, la etiqueta
  \textbf{no prueba que la transmitió ese dispositivo en particular}: en principio
  cualquiera que conozca $K$ (incluido otro dispositivo legítimo, o el propio colector)
  podría haber generado esa etiqueta. Con firma digital (parte a) cada dispositivo es el
  único que puede producir su firma (clave privada propia); con MAC compartido esa
  propiedad se pierde por completo.
  \item El colector mismo pasa a ser un blanco crítico: si se compromete, se filtra $K$ y
  cae la seguridad de \emph{todos} los dispositivos de la red, no solo de uno.
\end{itemize}

\textbf{b) (repetida) --- Evitar el replay de señales viejas.} El ataque descripto es un
\textbf{replay}: alguien corta una sección y coloca un aparato que retransmite las
firmas/etiquetas \textbf{viejas} capturadas de los dispositivos que quedaron desconectados,
haciendo que el colector siga viendo ``todo ok'' aunque esos dispositivos ya no estén. La
defensa (Clase 5) es agregar \textbf{frescura} (nonce o timestamp/contador) a lo que se
firma o etiqueta. Como el ancho de banda está fijado en 320 bits y no se puede agrandar, no
alcanza con agregar el contador ``en claro'' junto a la firma (ocuparía más bits); la
solución es la misma idea que justifica firmar el \textbf{hash} en la Clase 4: en vez de
firmar/etiquetar $\text{serie}_i$ solo, cada dispositivo firma/etiqueta
$H(\text{serie}_i\,\|\,c_i)$, donde $c_i$ es un \textbf{contador} (o ventana de tiempo)
que el dispositivo incrementa en cada transmisión y que el colector puede predecir
(ambos, dispositivo y colector, mantienen el contador sincronizado, o el colector acepta
una ventana de contadores cercana al último visto). Como el hash tiene salida de tamaño
fijo, la firma/etiqueta sigue ocupando exactamente 320 bits, sin importar que ahora el
mensaje firmado sea más largo. El colector, al verificar, recalcula
$H(\text{serie}_i\|c_i^{\text{esperado}})$ y ejecuta $\mathrm{Vrfy}$ (o recomputa el MAC)
contra ese valor; si el aparato del atacante repite una firma/etiqueta vieja, el contador
que usó para generarla ya \textbf{no coincide} con el esperado por el colector, y la
verificación falla (detecta el replay). Si el contador nunca avanza mientras se supone que
el dispositivo sigue transmitiendo, el colector puede además usar eso como señal adicional
de anomalía.

% =====================================================================
\subsection{Ejercicio --- Organizame (pentesting de una app de calendario)}
% =====================================================================

\begin{consigna}
``Organizame'' es una aplicación móvil que busca potenciar el uso eficiente del tiempo de
las personas.

La aplicación permite que un usuario se registre utilizando su cuenta de Google o su Apple
Id, y pide acceso al calendario de dicha cuenta para obtener información de las reuniones
agendadas.

La información de agenda es almacenada en una base de datos en el servidor, donde procesos
corren regularmente análisis y proponen mejoras a la agenda. Cuando uno de estos procesos
encuentra una posible mejora, envía una notificación a la aplicación para que el usuario
decida si desea aplicar el cambio o no. Algunos procesos son simples y rápidos, aunque hay
algunos que tienen un orden de complejidad $O(n^3)$ en la cantidad de eventos del mes.

La aplicación permite compartir huecos de tiempo libre del calendario de forma publica,
para que otras personas puedan agendar reuniones. Para ello disponibiliza una pagina web
por usuario, de la forma: \texttt{https://organizame.io/\{user\_id\}/public/calendar}. Un
usuario puede habilitar o deshabilitar dicha página en su configuración de perfil. En ella,
cualquier persona puede agendar una reunión en los momentos disponibles, y dejarle al dueño
un mensaje que verá en su calendario.

Ante un pedido de realizar un pentest que abarque el sistema, establecer 4 hipótesis de
falla \textbf{pertinentes} para el caso, basadas en problemas de seguridad en aplicaciones.
Explicar \textbf{para cada una}, como probarla según la metodología (la acción concreta que
ejecutaría para verificar o refutar la hipótesis). (4 puntos, 1 pto cada uno)

\textbf{Aclaración}: No utilizar problemas genéricos como hipótesis de falla. Plantear
situaciones y lugares concretos donde podría estar el problema. Por ejemplo, en lugar de
decir ``puede haber un buffer overflow'', decir ``es probable que haya un buffer overflow
en la lectura del parámetro x que debe interpretarse como número y puede tener un buffer
pequeño asumiendo que el número tiene menos de los dígitos de un long''.
\end{consigna}

\paragraph{Temas.}
Clase 6 (Segunda Parte) --- Pentesting (\S Metodología de hipótesis de falla, en
particular Paso 2 ``Hipótesis'' y Paso 3 ``Prueba''); Clase 1 --- Políticas y Control de
Acceso (\S OAuth 2.0 y \emph{scope}); Clase 3 --- Principios de Diseño y Vulnerabilidades
(\S Tipos de vulnerabilidades: relaciones de confianza, injection flaws; \S disponibilidad
y DoS como ataque de seguridad).

\paragraph{Qué necesitás saber.}
\begin{itemize}
  \item \textbf{Metodología de hipótesis de falla (memorizar):} Recolección de información
  $\to$ Planteo de hipótesis de falla $\to$ Prueba de la vulnerabilidad $\to$
  Generalización $\to$ Eliminación (opcional). Acá se pide directamente el paso 2
  (hipótesis) justificado con el paso 3 (cómo probarla): ``definir cómo probar'' prioriza
  \emph{analizar documentación o comportamiento} antes que intentar explotar, y la prueba
  debe ser \textbf{repetible y concluyente}.
  \item \textbf{Relaciones de confianza (Clase 3, tipos de vulnerabilidades):} ``no
  chequear autorización porque \emph{nadie va a adivinar la URL}'' es exactamente el
  patrón de la URL \texttt{/\{user\_id\}/public/calendar}: hay que verificar que el
  servidor efectivamente valide el flag de ``habilitado'' del usuario y no simplemente
  confíe en que el \texttt{user\_id} es difícil de adivinar.
  \item \textbf{OAuth y \emph{scope} (Clase 1):} el token de acceso debe limitarse al
  \emph{scope} mostrado en la pantalla de consentimiento (acá, lectura de reuniones);
  pedir de más (por ejemplo, permiso de escritura/administración total del calendario) es
  una falla de diseño de autorización.
  \item \textbf{Disponibilidad como propiedad de seguridad (Clase 3):} un ataque que tira
  o degrada el sistema (DoS) \emph{es} un ataque de seguridad aunque no comprometa
  confidencialidad ni integridad; un proceso de complejidad $O(n^3)$ sobre un dato que un
  atacante puede inflar (cantidad de eventos del mes) es candidato directo.
  \item \textbf{Injection flaws (Clase 3):} el profesor remarca que $\sim$90\% de las
  vulnerabilidades caen en esta familia; cualquier campo de entrada no controlado (acá, el
  ``mensaje'' que deja un desconocido en el calendario público) es un lugar concreto donde
  buscarla.
\end{itemize}

\paragraph{Notas y conexiones.}
Las cuatro hipótesis se ubican, cada una, en un componente concreto descripto en el
enunciado (la URL pública, el proceso $O(n^3)$, el flujo OAuth con Google/Apple, el campo
de mensaje libre), tal como pide la aclaración de no usar hipótesis genéricas. Todas están
tomadas de puntos ya vistos en el material (relaciones de confianza, disponibilidad/DoS,
scope de OAuth, injection flaws) y no se agrega nada fuera de eso.

\paragraph{Resolución.}

\textbf{Hipótesis 1 --- IDOR / relación de confianza en \texttt{/\{user\_id\}/public/calendar}.}
Es probable que el servidor, al recibir un \texttt{GET} a esa URL, no vuelva a validar que
el usuario efectivamente tenga habilitada la opción ``pública'' en su perfil, confiando en
que \texttt{user\_id} es difícil de adivinar o de enumerar (la relación de confianza
``nadie va a adivinar la URL'' que señala el material). Si los \texttt{user\_id} son
secuenciales o de algún modo predecibles, cualquiera podría ver/agendar en calendarios que
el dueño nunca hizo públicos.
\emph{Cómo probarla:} crear una cuenta de prueba con la página pública \textbf{deshabilitada},
y con otra cuenta (o sin sesión) acceder directamente a
\texttt{https://organizame.io/\{ese\_user\_id\}/public/calendar}; si igual devuelve datos
de agenda o permite agendar, queda demostrado. También probar con \texttt{user\_id}
consecutivos a uno propio habilitado, para ver si expone calendarios de terceros que no lo
sabían.

\textbf{Hipótesis 2 --- Denegación de servicio vía el proceso $O(n^3)$.} Es probable que el
proceso de análisis de complejidad $O(n^3)$ en la cantidad de eventos del mes no tenga un
límite de eventos por usuario/mes, por lo que un usuario (o un atacante que agenda a través
de la página pública de un tercero) podría inflar artificialmente la cantidad de eventos de
un mes para que ese proceso batch consuma recursos desproporcionados, degradando la
disponibilidad del servicio para otros usuarios.
\emph{Cómo probarla:} con una cuenta de prueba, generar mediante un script una gran
cantidad de eventos en un mismo mes (por API propia o a través del flujo de agendar
reunión desde una página pública), y medir el tiempo/CPU que insume la corrida del proceso
de análisis para ese mes, verificando si crece de forma cúbica y si llega a afectar al
resto del sistema (colas, timeouts, caída del servicio).

\textbf{Hipótesis 3 --- Scope de OAuth excesivo al integrar Google/Apple Calendar.} Es
probable que, al pedir ``acceso al calendario'' de la cuenta de Google/Apple, la app
solicite un \emph{scope} más amplio del necesario (por ejemplo, lectura \textbf{y}
escritura/administración del calendario, o acceso a otros datos de la cuenta), cuando para
leer reuniones agendadas alcanzaría con un scope de solo lectura.
\emph{Cómo probarla:} iniciar el flujo de registro con Google/Apple y revisar la pantalla
de consentimiento y/o decodificar el \emph{access token} obtenido, verificando qué scopes
concretos se solicitan y comparándolos con lo mínimo necesario (leer eventos); si el scope
permite modificar o borrar eventos de la cuenta real de Google/Apple del usuario sin que
eso sea necesario para la funcionalidad descripta, queda demostrado el exceso.

\textbf{Hipótesis 4 --- Injection en el campo de mensaje de la página pública.} Es probable
que el campo de texto donde un desconocido deja un mensaje al agendar una reunión desde
\texttt{/\{user\_id\}/public/calendar} no sanitice la entrada antes de guardarla o de
mostrarla en el calendario del dueño, permitiendo una inyección (por ejemplo, un payload de
script si el mensaje se renderiza en la vista web del calendario del dueño).
\emph{Cómo probarla:} desde la página pública, agendar una reunión de prueba dejando como
mensaje un payload típico de prueba (por ejemplo una etiqueta de script inocua, o una
comilla simple si el mensaje se guarda en una base de datos SQL), y luego revisar, desde la
cuenta del dueño, si el mensaje se ejecuta/se interpreta en vez de mostrarse como texto
plano, o si genera un error de base de datos que delate una inyección.

\begin{ojo}
No se incluyen hipótesis genéricas (``puede haber XSS en algún lado'', ``puede haber SQL
injection en algún formulario'') porque la consigna explícitamente lo prohíbe: cada
hipótesis de arriba está anclada a un componente y un flujo concreto descripto en el
enunciado de ``Organizame''.
\end{ojo}
